\documentclass[11pt, a4paper]{article}

\usepackage[T1]{fontenc}
\usepackage[utf8]{inputenc}
\usepackage{tgpagella}
\usepackage[spanish, es-noshorthands]{babel}
\usepackage{amsmath, amssymb, bm}
\usepackage{tabularx}
\usepackage[table, xcdraw]{xcolor}
\usepackage[most]{tcolorbox}
\usepackage[margin=2.5cm]{geometry}
\usepackage{ragged2e}
\usepackage{fancyhdr}

\pagestyle{fancy}
\fancyhf{}
\setlength{\headheight}{16pt}
\lhead{\textbf{UCB Sede Tarija} | Ing. Mecatrónica}
\chead{}
\rhead{IMT-121 Circuitos Electrónicos I | Pre-Brief Exp. B}
\lfoot{Prof: M.Sc. Ing. Bernardo Quiroga Turdera}
\rfoot{Página \thepage}
\renewcommand{\headrulewidth}{0.4pt}

\definecolor{ucbblue}{RGB}{0, 51, 102}

\begin{document}

\begin{tcolorbox}[colback=ucbblue!5!white, colframe=ucbblue, title=\textbf{Prelaboratorio Conceptual: Experiencia Integrada B}]
    \centering \textbf{Preparación Analítica para Máxima Transferencia de Potencia}
\end{tcolorbox}

\section*{1. El Puente desde la Experiencia A}
Con la Experiencia A, su equipo logró reducir una red compleja de dos fuentes a un simple equivalente de Thévenin: un voltaje $\bm{V_{th}}$ en serie con una resistencia $\bm{R_{th}}$.

Si conectamos una carga variable $R_L$ a este equivalente, la corriente de carga es $I_L = \frac{V_{th}}{R_{th} + R_L}$ y el voltaje es $V_L = I_L R_L$. La potencia eléctrica útil entregada a esta carga es:
\begin{equation*}
    P_L = I_L^2 R_L = V_{th}^2 \left( \frac{R_L}{(R_{th} + R_L)^2} \right)
\end{equation*}

\section*{2. Derivación del Máximo (Demostración Obligatoria)}
Para hallar en qué valor de carga se extrae la potencia máxima, debemos encontrar el punto donde la derivada de la potencia respecto a $R_L$ es cero ($\frac{dP_L}{dR_L} = 0$).

Aplicando la regla del cociente:
\begin{align*}
    \frac{dP_L}{dR_L} &= V_{th}^2 \left[ \frac{1 \cdot (R_{th} + R_L)^2 - R_L \cdot 2(R_{th} + R_L)}{(R_{th} + R_L)^4} \right] \\
    &= V_{th}^2 \left[ \frac{(R_{th} + R_L) - 2R_L}{(R_{th} + R_L)^3} \right] \\
    &= V_{th}^2 \left[ \frac{R_{th} - R_L}{(R_{th} + R_L)^3} \right] = 0
\end{align*}
Para que esta expresión sea cero (y dado que $V_{th} \neq 0$ y la resistencia es finita), el numerador debe ser cero:
\begin{equation*}
    R_{th} - R_L = 0 \implies \mathbf{R_L = R_{th}}
\end{equation*}

\section*{3. Preparación para la Próxima Sesión}
Responda a las siguientes preguntas con su equipo y anótelas en su cuaderno. Sin esta planificación, no podrán iniciar la Experiencia B.

\begin{enumerate}
    \item Utilizando el $\bm{R_{th}}$ que acaban de calcular y medir en la Experiencia A, ¿qué valor exacto de carga $R_L$ predicen que extraerá la máxima potencia?
    \item Para la siguiente sesión necesitaremos "barrer" resistencias para observar la curva. Seleccionen de su kit \textbf{5 valores comerciales} de resistencias:
    \begin{itemize}
        \item Dos valores claramente por debajo de su $R_{th}$ (Régimen $R_L \ll R_{th}$).
        \item Un valor lo más exacto posible a su $R_{th}$ (Régimen Adaptado).
        \item Dos valores claramente por encima de su $R_{th}$ (Régimen $R_L \gg R_{th}$).
    \end{itemize}
    \item Diseñe la \textbf{Tabla Mínima de Laboratorio} que usarán la próxima clase. Debe contener columnas para: $R_L$ (nominal), $R_L$ (medida), $V_L$, $I_L$, $P_L$ y Eficiencia $\eta$.
\end{enumerate}

\end{document}
